% PLOS ONE manuscript -- PONE-D-26-28965 (revision)
% Based on the official PLOS LaTeX template v3.8 (Apr 2026).
%
% % % % % % % % % % % % % % % % % % % % % % %
%
% THIS IS THE CLEAN MANUSCRIPT.
%
% It contains no revision highlighting: PLOS requires the clean "Manuscript"
% file to be free of tracked changes.
%
% Graphics ARE embedded here, exactly as in initial_submission/main.tex, so the
% working PDF is readable. PLOS requires each figure to be uploaded as an
% individual file (https://journals.plos.org/plosone/s/figures), that project is
% generated by "make prepare-plos-submission", which strips \includegraphics and
% writes figures/Fig<N>.tif. Do not strip the graphics by hand.
%
% The "Revised Manuscript with Track Changes" file is produced from this source
% with latexdiff against the previously submitted version. See ../README.md.
% Do not hand-edit colors into this file.
%
% % % % % % % % % % % % % % % % % % % % % % %
%
% NOTE ON \rev / \revs
% The wrappers \rev{...} and \revs{...} mark passages that were added or
% modified in the first revision round. They are deliberately defined as
% identity macros so that they never render color, while preserving the record
% of which passages changed. Keep them. Do not re-enable the colored versions
% in this file.
%
% % % % % % % % % % % % % % % % % % % % % % %

\documentclass[10pt,letterpaper]{article}
\usepackage[top=0.85in,left=2.75in,footskip=0.75in]{geometry}

% amsmath and amssymb packages, useful for mathematical formulas and symbols
\usepackage{amsmath,amssymb}

% Use adjustwidth environment to exceed column width (see example table in text)
\usepackage{changepage}

% textcomp package and marvosym package for additional characters
\usepackage{textcomp,marvosym}

% cite package, to clean up citations in the main text. Do not remove.
\usepackage{cite}

% Use nameref to cite supporting information files (see Supporting Information section for more info)
\usepackage{nameref,hyperref}

% line numbers
\usepackage[right]{lineno}

% ligatures disabled
\usepackage[nopatch=eqnum]{microtype}
\DisableLigatures[f]{encoding = *, family = * }

% color can be used to apply background shading to table cells only
\usepackage[table]{xcolor}

% array package and table rules
\usepackage{array,booktabs,longtable,float,placeins}

% Round-1 revision markers, defined as identity macros so that the clean
% manuscript contains no in-text color (PLOS ONE requirement for the
% unmarked "Manuscript" file). \rev{...} wraps changed paragraphs and
% \revs{...} wraps changed inline fragments such as headings.
\newcommand{\rev}[1]{#1}
\newcommand{\revs}[1]{#1}

% multirow, for scenario cells spanning several rows in the results table
\usepackage{multirow}

% create "+" rule type for thick vertical lines
\newcolumntype{+}{!{\vrule width 2pt}}

% left-aligned fixed-width column, used by the tables in tables/
\newcolumntype{L}[1]{>{\raggedright\arraybackslash}p{#1}}

% create \thickcline for thick horizontal lines of variable length
\newlength\savedwidth
\newcommand\thickcline[1]{%
  \noalign{\global\savedwidth\arrayrulewidth\global\arrayrulewidth 2pt}%
  \cline{#1}%
  \noalign{\vskip\arrayrulewidth}%
  \noalign{\global\arrayrulewidth\savedwidth}%
}

% \thickhline command for thick horizontal lines that span the table
\newcommand\thickhline{\noalign{\global\savedwidth\arrayrulewidth\global\arrayrulewidth 2pt}%
\hline
\noalign{\global\arrayrulewidth\savedwidth}}


% PLOS ONE requires double-spaced manuscript text. The two lines below may be
% commented out while working to keep the PDF compact. The tracked-changes route
% ("make prepare-marked-manuscript") uncomments them again.
\usepackage{setspace}
\doublespacing

% Text layout
\raggedright
\setlength{\parindent}{0.5cm}
\textwidth 5.25in
\textheight 8.75in

% Bold the 'Figure #' in the caption and separate it from the title/caption with a period
% Captions will be left justified
\usepackage[aboveskip=1pt,labelfont=bf,labelsep=period,justification=raggedright,singlelinecheck=off]{caption}
\renewcommand{\figurename}{Fig}

% Supporting information items are uploaded as separate files, so PLOS labels
% them "S1 Table." rather than continuing the main table numbering. The
% submission package replaces them with legend paragraphs carrying that label
% (make prepare-plos-submission); this switch makes the working PDF and the
% tracked-changes file, which still render the items inline, number them the
% same way, so the in-text citations read alike in all three.
\DeclareCaptionLabelFormat{plossupporting}{S#2 Table}
\newcommand{\supportinginformationstyle}{%
  \captionsetup[table]{labelformat=plossupporting,labelsep=period}%
  \setcounter{table}{0}%
  \renewcommand{\thetable}{\arabic{table}}%
}

% Please use the included `plos2025.bst` as your BibTeX style
\bibliographystyle{plos2025}

% Remove brackets from numbering in List of References
\makeatletter
\renewcommand{\@biblabel}[1]{\quad#1.}
\makeatother



% Header and Footer with logo
\usepackage{lastpage,fancyhdr,graphicx}
\usepackage{epstopdf}
%\pagestyle{myheadings}
\pagestyle{fancy}
\fancyhf{}
%\setlength{\headheight}{27.023pt}
%\lhead{\includegraphics[width=2.0in]{PLOS-submission.eps}}
\rfoot{\thepage/\pageref{LastPage}}
\renewcommand{\headrulewidth}{0pt}
\renewcommand{\footrule}{\hrule height 2pt \vspace{2mm}}
\fancyheadoffset[L]{2.25in}
\fancyfootoffset[L]{2.25in}
\lfoot{\today}

%% Include all user-defined macros below

%% END MACROS SECTION


\begin{document}
\vspace*{0.2in}

% Title must be 250 characters or less.
\begin{flushleft}
{\Large
\textbf{Zero-shot generalization of EEG attention-state decoding across increasingly stringent validation strategies on open datasets} % Please use "sentence case" for title and headings.
}
\newline
% Insert author names, affiliations and corresponding author email (do not include titles, positions, or degrees).
\\
Sanzhar Sovet\textsuperscript{1},
Beibit Abdikenov\textsuperscript{1*},
Manzura Zholdasova\textsuperscript{2,3},
Altyngul Kamzanova\textsuperscript{2},
Medet Mukushev\textsuperscript{1},
Almira Kustubayeva\textsuperscript{2,3}
\\
\bigskip
\textbf{1} Science and Innovation Center ``Artificial Intelligence'', Astana IT University, Astana, Kazakhstan
\\
\textbf{2} Department of Biophysics, Biomedicine, and Neuroscience, Al-Farabi Kazakh National University, Almaty, Kazakhstan
\\
\textbf{3} Brain Institute, Al-Farabi Kazakh National University, Almaty, Kazakhstan
\\
\bigskip

% Use the asterisk to denote corresponding authorship and provide email address in note below.
* beibit.abdikenov@astanait.edu.kz

\end{flushleft}
% Please keep the abstract below 300 words

\section*{Abstract}

Electroencephalography (EEG) decoding models must generalize across people and experimental paradigms, yet non-stationary signals and the uncertain equivalence of labels such as ``attention'' make this difficult. Evidence on whether nominally aligned attention-related datasets provide useful zero-shot transfer across tasks and datasets remains limited. Here, we evaluated zero-shot generalization of an operationally defined binary contrast between high- and low-attention states in two technically compatible public EEG datasets: a five-participant mental attention-state dataset recorded during a train-simulator task and SAM-40, a 40-participant dataset spanning Stroop, arithmetic, and mirror-image tasks. The datasets were harmonized through a shared 12-channel input space, a fixed preprocessing pipeline, and a common label mapping. Drowsiness was excluded to reduce confounding by arousal, visual input, and ocular activity. Five classifiers spanning linear, tree-based, and deep-learning approaches were assessed using group-disjoint within-dataset, cross-subject, cross-task, and cross-dataset validation, without target-specific adaptation. Mean balanced accuracy across decoders was nearly unchanged when generalizing to unseen participants, differing from the within-dataset baseline by at most two percentage points. Compared with target-specific baselines, balanced accuracy fell by an average of 6.4 percentage points between SAM-40 task compositions and 13.7 points across datasets; every directed transfer comparison showed a decline. Absolute accuracy also varied materially across decoders. Thus, shared channels, technically compatible acquisition, and nominally aligned labels were insufficient to preserve target-specific performance under heterogeneous transfer. These directed pre-adaptation benchmarks support empirical source selection and provide matched references for evaluating whether transfer learning, domain adaptation, or personalization yields meaningful gains. The findings concern protocol-defined contrasts rather than a unitary neural marker of attention.

\clearpage
\newgeometry{top=0.85in,left=1in,right=1in,footskip=0.75in}
\linenumbers

% Use "Eq" instead of "Equation" for equation citations.
\section*{Introduction}

Electroencephalography (EEG) is a technique for recording brain activity using multichannel scalp electrodes that measure voltage signals reflecting underlying cognitive processes. It is widely used in neuroscience and brain--computer interface (BCI) research because of its non-invasive nature, high temporal resolution, portability, and relatively low cost~\cite{rashid2020bci}. At the same time, EEG signals are inherently noisy, non-stationary, and exhibit high inter-subject and inter-session variability~\cite{wu2022transfer,siino2025rdwt}. In addition, these challenges are compounded by the scarcity of large, diverse, and well-annotated open EEG datasets~\cite{sun2025attention}. Consequently, developing a robust decoding model with strong generalization capabilities across subjects, sessions, and tasks remains an open and central challenge~\cite{li2025foundation,wu2022transfer}. One promising approach to address these problems is transfer learning (TL). TL reuses representations already learned on one or more source problems to reduce the amount of labeled data and training effort a related target problem would otherwise require~\cite{zhang2020transfer}. Its success depends not only on modeling and learning strategies but also on the degree of domain and label compatibility between source and target datasets.

For cognitive constructs such as attention, this compatibility is especially consequential and difficult to guarantee because attention remains an ambiguous concept and is not a unitary construct with a single universally accepted operational definition~\cite{sun2025attention,hommel2019noone}. Rather, attention is an umbrella term encompassing a set of partially independent processes and functions, such as selecting relevant information, maintaining focus over time, shifting attention between tasks, and allocating cognitive resources. These processes may operate on both external sensory input and internally represented information, including information maintained in working memory~\cite{hommel2019noone,knudsen2007components,petersen2012attention}. The relative contributions of these processes may vary with task demands, sensory modalities, task goals, experimental instructions, assigned conditions, self-reported assessments, and other aspects of experimental design and label definition. Consequently, experimental paradigms described as measuring ``attention'' may probe non-equivalent configurations of cognitive processes. As a result, different datasets may represent different attention-related cognitive states while assigning them the same nominal labels. \rev{In this study, we therefore treat ``high-attention'' and ``low-attention'' as operationally defined protocol conditions rather than as an independently validated attention construct, and ask whether decoders discriminate these conditions rather than whether they decode attention itself.}

Previous EEG transfer-learning research has been predominantly concerned with cross-subject and cross-session generalization, whereas cross-task, cross-dataset, and cross-paradigm transfer have received comparatively less systematic attention~\cite{wu2022transfer,zhou2022crosstask,song2024signal}. Other recent work has evaluated cross-subject robustness of deep decoders on multiple motor-imagery benchmark datasets while keeping the task paradigm fixed~\cite{siino2025ratiowavenet}. Consequently, systematic empirical evidence for source--target suitability across heterogeneous attention-related EEG datasets and tasks remains limited. Source--target suitability is often motivated conceptually or inferred from nominally aligned labels. This leaves unclear which source--target combinations provide useful zero-shot transfer, how much target discriminability is retained relative to target-specific baselines, and whether transferability is direction-dependent. Establishing such pre-adaptation baselines is necessary to justify source selection and to interpret gains or negative transfer from subsequent transfer-learning and domain-adaptation methods~\cite{zhang2023negative}.


Here, we address this gap by systematically evaluating zero-shot generalization across two technically compatible public EEG datasets spanning multiple tasks. \rev{Under a common preprocessing pipeline, we place within-dataset baseline performance, cross-subject generalization, and zero-shot transfer across tasks and datasets on a single ladder of increasingly stringent validation strategies. We characterize how performance behaves along this ladder, comparing directed source--target pairs against target-specific baselines to establish whether transfer is direction-dependent, and we test whether the resulting pattern is reproduced across five decoders ranging from linear and tree-based baselines to deep neural networks.} Together, these analyses yield reproducible pre-adaptation benchmarks for evaluating candidate sources and contextualizing both gains and negative transfer produced by subsequent transfer-learning and domain-adaptation methods.


\section*{Open-access datasets}

\subsection*{EEG data for mental attention state detection (Dataset A)}

The dataset described in~\cite{aci2019attention} comprises 34 EEG recordings collected from five participants over the course of seven days, with the fifth participant completing only six sessions. During each session, participants operated a virtual train using the \textit{Microsoft Train Simulator}. Each recording lasted approximately 35 to 55 minutes. The participant cohort consisted of healthy university student volunteers who provided written informed consent after receiving information about the study objectives and procedures; the original study reported that it was conducted in accordance with the ethical guidelines of Mersin University~\cite{aci2019attention}.

During every session, participants followed the same predefined route in the virtual train simulator. The route required little active control and kept task demands comparable across recordings. Each session comprised three sequential phases designed to reflect distinct target mental states: \textit{focused}, \textit{unfocused}, and \textit{drowsy}. Participants received standardized instructions from the experiment supervisor for each phase:

\begin{itemize}
  \item \textbf{Focused (0--10 min):} The focused phase required active monitoring of and mental engagement with the simulator, including sustained attention to the virtual controls and screen events.
  \item \textbf{Unfocused (10--20 min):} During the unfocused phase, participants disengaged from the task, ceased interaction, and were instructed to keep their eyes open and remain awake.
  \item \textbf{Drowsy (20--30 min):} During the drowsy phase, participants were allowed to relax, close their eyes, and enter a drowsy or lightly dozing state, simulating reduced cognitive arousal.
\end{itemize}

EEG signals were recorded using a modified Emotiv headset, equipped with wet electrodes and supporting wireless Bluetooth transmission. The device recorded data at a sampling rate of 128~Hz and with a voltage resolution of 0.51~$\mu$V. The dataset includes 14 EEG channels, labeled according to the 10--20 international system.

\subsection*{SAM-40: dataset of 40-subject EEG recordings (Dataset B)}

The \textbf{SAM-40} dataset, described in~\cite{ghosh2022sam40}, contains EEG recordings from 40 healthy participants (14 females and 26 males) with a mean age of 21.5 years. The original study obtained informed consent from all participants and approval from the Institutional Ethics Committee of Gauhati University (GUIEC/2019/019, 15 October 2019). It reported voluntary and anonymous participation and compliance with the Declaration of Helsinki~\cite{ghosh2022sam40}. The dataset was collected to investigate short-term stress responses induced by various cognitive tasks. Its protocol corresponded to the \textit{continuous cognitive engagement} paradigm described in~\cite{chen2025hybrid} and included four conditions:

\begin{itemize}
  \item \textbf{Resting state}: Participants relaxed while listening to music during EEG recording.
  \item \textbf{Stroop Color-Word Test}: Participants identified the color of the ink in which a word was printed while ignoring the word's semantic meaning. The task included two types of stimuli: \textit{congruent} (the color name matched the ink color) and \textit{incongruent} (the color name and ink color differed). Each task run included 11 stimuli.
  \item \textbf{Arithmetic problem solving}: Participants mentally solved a mathematical problem and determined whether the presented answer was correct, responding with a ``thumbs up'' or ``thumbs down'' gesture. Each task run included 6 arithmetic stimuli.
  \item \textbf{Mirror Image Recognition}: Participants were shown images and asked to determine whether they were mirror-symmetric or asymmetric, responding with the same gestures. Each task run included 8 images.
\end{itemize}

Each task was performed for 25 seconds. Between consecutive cognitive tasks, participants rested for 5 seconds and then received 10 seconds of instructions for the next task. After each trial, participants rated the stress experienced during each task on a 10-point scale, where 1 indicated minimal stress and 10 indicated maximal stress. Three trials, each using a different set of stimuli, were recorded for every participant.

EEG signals were recorded using an Emotiv Epoc Flex headset with 32 wet electrodes, labeled according to the international 10--20 system. Data were acquired at a sampling rate of 128~Hz and a voltage resolution of 0.51~$\mu$V.

\subsection*{Ethics statement}

This study was a secondary analysis of publicly available EEG recordings and involved no additional data collection or participant contact. No additional ethical approval or participant consent was required for this secondary analysis. The consent procedures and ethical oversight reported for the original studies are described in the mental attention-state dataset and SAM-40 sections above.

\section*{Methodology}

The methodological design of this study was developed to evaluate the generalizability and transferability of EEG-based attention decoding across different recording conditions, participants, and task paradigms. To this end, we conducted a secondary analysis of two publicly available datasets that differ in their experimental protocols but share compatible technical specifications, enabling a controlled comparative analysis.

\rev{The datasets were harmonized through a binary operational labeling scheme that
distinguishes high-attention states, defined by protocol-imposed task demand,
from low-attention states of rest or task disengagement, enabling consistent
evaluation across scenarios. These labels denote protocol-defined conditions
rather than independently validated levels of attention. They may reflect
differences in task demands, arousal, sensory input, stress, or motor responses
between resting and active cognitive conditions. The analyses therefore assess
whether decoders discriminate the specified protocol conditions, rather than
whether they decode attention as a general psychological construct.}

\rev{To evaluate zero-shot generalization, we used the cross-subject,
cross-task, and cross-dataset validation protocols. In each protocol, a
decoder was trained on its source composition and applied unchanged to the
held-out target data. The within-domain baseline was target-specific and
served as a reference for these zero-shot evaluations. Held-out target
recordings contributed to only one fitted preprocessing operation, ICA artifact
cleaning. ICA was fitted without class labels, separately for each selected
session in the mental attention-state dataset and for the selected recordings
of each participant in SAM-40. No ICA decomposition pooled source and target
recordings or recordings from different participants. Feature extraction was
applied in the same form to source and target recordings. All other fitted
operations were estimated from source training data only, including
standardization statistics, class weights, and decoder parameters. No target
data were used for model selection, hyperparameter search, early stopping,
checkpoint selection, decoder retraining, or tuning. Target labels were never
used. Thus, the zero-shot protocols tested direct application of a
source-trained decoder without target-specific decoder adaptation.}

All experiments were executed through a version-controlled computational workflow. For every run, the workflow stores the resolved configuration of data preparation, validation, and model settings together with the run outputs needed to reconstruct the reported evaluation metrics. These records provide run-level provenance and allow each reported metric to be traced to the experiment that produced it. The complete implementation, experiment configurations, and run-level artifacts are available at \url{https://github.com/sovetskiysn/eeg-validation-strategies}, released under the MIT license and archived at \url{https://doi.org/10.5281/zenodo.22255472}.

\subsection*{Dataset alignment and integration rationale}

Effective evaluation of cross-task and cross-dataset generalization relies on the availability of datasets that are compatible both technically and conceptually. In this study, we used two publicly available EEG datasets that differ in their original purpose and labeling schemes, yet retain sufficient technical and conceptual compatibility for a meaningful comparison.

A major strength of this dataset combination is the compatible acquisition setup of the two datasets. Both datasets were recorded using Emotiv EEG systems with the same sampling rate (128~Hz) and voltage resolution, although their original electrode montages differed (14 versus 32 channels). For cross-dataset analyses, only the 12 channels common to both datasets were retained. Differences in EEG recording systems and electrode configurations can hinder cross-dataset transferability~\cite{lan2019domain}. Restricting the analysis to the shared channels substantially reduces this configuration-related mismatch, although it does not eliminate all acquisition-related distribution shifts. Within the limited set of publicly available datasets suitable for this comparison, this pair offered a comparatively compatible basis for the present comparative evaluation.

The conceptual rationale for combining the datasets rests on the cognitive conditions represented in their respective protocols. The mental attention-state dataset contains explicit attention-state labels (\textit{Focused}, \textit{Unfocused}, and \textit{Drowsy}) collected during a low-interaction simulation task, whereas SAM-40 was originally designed for stress assessment. Both include conditions involving attentional processes, although their original constructs differ. The SAM-40 dataset includes tasks widely recognized in cognitive neuroscience as requiring attentional engagement:
\begin{itemize}
    \item \textbf{Stroop Color-Word Test}: probes both \textit{selective} and \textit{sustained attention}, requiring participants to suppress automatic reading responses and continually focus on the ink color throughout the task~\cite{macleod1991stroop, scarpina2017stroop, esterman2013zone}.
    \item \textbf{Mental arithmetic task}: demands \textit{sustained internal attention} and working memory for continuous manipulation of numerical information~\cite{grabner2007math, kitaura2017functional}.
    \item \textbf{Mirror Image Recognition}: engages visuospatial processing and requires both \textit{selective} and \textit{sustained attention} to discriminate between symmetric and asymmetric patterns, relying on mental rotation processes that demand continuous cognitive effort~\cite{shepard1971mental, hyun2007visual, tomasino2016mental}.
\end{itemize}

The simulator-based and task-based datasets correspond, respectively, to the \textit{naturalistic stimulus-based protocols} and \textit{continuous cognitive engagement} paradigms described in~\cite{chen2025hybrid}. In the unified labeling scheme, \textit{Rest} and \textit{Unfocused} are treated as low-attention states, whereas \textit{Focused}, \textit{Stroop}, \textit{Arithmetic}, and \textit{Mirror recognition} are treated as high-attention states. The \textit{Drowsy} phase is excluded from this mapping and from all subsequent preprocessing and analyses because permitted eye closure and drowsiness may confound the comparison with differences in arousal, visual input, and ocular activity. This mapping is summarized in Table~\ref{tab:attention_label_mapping}. \rev{Throughout this paper, ``high-attention'' and ``low-attention'' states denote these operationally defined protocol conditions rather than an independently validated attention construct.}

\input{tables/attention_label_mapping}
\FloatBarrier

\subsection*{Data preparation}

\rev{The same deliberately simple preparation pipeline was applied to
both datasets and across all validation protocols. It used common filtering
settings, ICA fitting and component-exclusion criteria, channel-alignment
selection, and windowing parameters. These settings were fixed and were not tuned
separately for a dataset, task composition, validation protocol or decoder. The
purpose was to place recordings in a common input space while limiting
preprocessing itself as an additional source of variation. Differences between
validation results are therefore interpreted against a fixed preparation
procedure. The pipeline was instantiated as a fixed set of dataset
preparations, namely the dataset compositions evaluated in every validation
protocol and listed in
Table~\ref{tab:dataset_compositions}. The remainder of this subsection
describes the selection, standardization, cleaning and segmentation steps that
produced them.}

\input{tables/dataset_compositions}

\rev{\textbf{Recording selection and label harmonization.} Recordings were selected
before signal processing. Following the source study, the first two recording
sessions for each participant in the mental attention-state dataset were
intended for habituation to the simulator and were not included in the
analyses~\cite{aci2019attention}. The \textit{Drowsy} condition was excluded
from the binary classification analysis because participants were
permitted to close their eyes and enter a drowsy state. This condition differs
from the \textit{Focused} and \textit{Unfocused} conditions not only in
protocol-defined task demand, but also in arousal, visual input and ocular
activity. Including it in the low-attention class would therefore confound the
intended contrast. The drowsy segment was removed before filtering and ICA, so
no drowsy epochs were generated or presented to a classifier. In contrast, the complete collection of SAM-40 recordings was
retained for analysis. SAM-40 provides the task recordings of each trial as
separate files. Within each scenario, the task recordings selected from one
trial were grouped into a single recording unit, whereas task recordings not
selected by the scenario were not carried forward to later preparation steps.
Accordingly, the binary mapping in Table~\ref{tab:attention_label_mapping} was applied to the
retained recordings from the two datasets.}

\rev{\textbf{Signal standardization and channel alignment.} Before filtering
and ICA cleaning, each recording was restricted to the 12 channels shared by
the two datasets: F3, F4, F7, F8, FC5, FC6, O1, O2, P7, P8, T7 and T8. Both
source datasets were originally sampled at 128~Hz and used the same documented
conversion factor of 0.51~$\mu$V per digit. Therefore, no resampling or
additional amplitude scaling was required for cross-dataset harmonization.}

\rev{\textbf{Filtering and ICA cleaning.} Filtering was applied to each
recording, and ICA was fitted within participant-specific acquisition
groupings. For the mental attention-state dataset, each grouping corresponded
to a selected session, whereas for SAM-40 it comprised the selected task recordings
of a participant. Training and test data therefore
never shared a decomposition, except in SAM-40 baseline validation, where
different runs of the same participant could fall into training and test folds,
so this baseline may be favored in comparisons with participant-disjoint protocols. To prevent preprocessing from becoming an additional source
of variation, we used a deliberately restrained, conventional procedure:
standard filtering followed by ICA-based artifact cleaning. Keeping
preparation deliberately simple allows differences in performance to be
interpreted primarily in relation to the recordings and validation protocols,
rather than to a complex preparation pipeline. EEG channels were notch-filtered
at 50~Hz, band-pass filtered from 1 to 45~Hz and resampled to 128~Hz. Before
ICA, they were rereferenced to the average reference. Artifact correction used independent component analysis with
extended-Infomax optimization, retaining components that explained 99\% of the
variance. The resulting components were automatically classified from their
spatial, spectral and temporal characteristics. Components identified as
eye-blink or muscle artifacts with a probability of at least 0.80 were removed.
The ICA random state was fixed by the experiment-wide configuration.}

\rev{\textbf{Window segmentation and metadata.} After ICA cleaning, each
retained labeled block was segmented into 5-s windows with 0.5-s overlap.
Every retained window inherited the label of its source block. Metadata
identifying its dataset, participant, session, task, run, recording unit and
temporal onset were retained with each window. These metadata were used to
support protocol-specific grouping and split definitions. In particular, the
recording-unit identifier ensured that all overlapping windows from a unit
remained on the same side of a baseline split.}

\rev{The fixed selection, standardization, cleaning and segmentation steps therefore
yielded a deterministic analytical window set for each dataset composition.
All model-specific input representations were derived from these same prepared
windows, rather than from separately prepared data. The retained recording
units, window counts and observed class distributions of the resulting
compositions are those reported in Table~\ref{tab:dataset_compositions}. Class weights were calculated
independently within each training fold. All SAM-40 compositions contain the same 40 participants and
differ only in the task recordings retained alongside the resting-state
recordings.}

\subsection*{\revs{Validation strategies}}

\rev{We evaluated the decoders under four validation protocols, each
addressing a distinct generalization question. The within-domain baseline
held out recording units from the same dataset and condition composition;
cross-subject validation changed participant identity; cross-task validation
changed both participant identity and SAM-40 task composition; and
cross-dataset validation changed datasets, participant populations, and
condition definitions. The participant, session, task, run, and recording-unit
metadata retained with each window defined the required groups. All windows
from a recording unit remained in the same partition, preventing overlapping
windows from crossing the training--test boundary~\cite{chen2025hybrid}.
Because baseline folds could include other recordings of test participants,
comparisons with the baseline also include a shift to unseen participants.
Cross-task and cross-dataset differences therefore cannot be attributed to
task composition or dataset alone. The protocol-specific splits and ICA fits
are detailed below.}

\textbf{Baseline validation.}
\rev{Performance was assessed on held-out recording units within the same dataset and condition composition. Selected recordings were divided into five stratified, group-disjoint folds, separately for the mental attention-state dataset, the SAM-40 dataset, and each task-specific composition of SAM-40. In the mental attention-state dataset, each selected participant--session recording was one group. In SAM-40, each participant--run was one group containing the resting-state recording and all task data included in the selected composition for that run. For example, when evaluating the Stroop composition, the resting-state and Stroop recordings from the same participant and run were always assigned to the same fold. In each fold, approximately four fifths of the recording units were used for training and the remaining units for testing. All windows from a recording unit were assigned to the same partition. A participant could contribute different sessions or runs to both training and testing, but the same physical recording could never occur in both partitions. Training and testing therefore shared the dataset and condition composition while differing in the recording units presented to the decoder. For data-dependent preprocessing, each selected session in the mental attention-state dataset was cleaned using an ICA decomposition fitted separately for that session, so the decomposition remained on one side of the split. In SAM-40, the decomposition used for cleaning was fitted across each participant's retained runs within the evaluated composition and could therefore span training and test runs from that participant. Keeping each unit intact prevented temporally adjacent or partially overlapping windows from the same acquisition from being assigned to opposite folds, thereby avoiding window-level temporal leakage. These estimates serve as within-domain comparators for the transfer protocols.}

\textbf{Cross-subject validation.}
\rev{In each leave-one-subject-out fold, the decoder was trained on all selected data from the remaining participants and tested on all selected data from the held-out participant. Data-dependent preprocessing of the held-out participant's recordings was limited to ICA fitted on that participant's own unlabeled data, separately for each selected session in the mental attention-state dataset and for the selected recordings of each participant in SAM-40. No decomposition used to clean source data included the held-out participant. The resulting performance quantifies zero-shot generalization to an unseen participant.}

\textbf{Cross-task validation.}
\rev{In SAM-40, each composition represented the same binary contrast of
low-attention rest versus high-attention task conditions, rather than
classification of task identity. Cross-task validation assessed directed
transfer between different compositions of these conditions. The analysis
included transfers between the Stroop, Arithmetic, and Mirror task-specific
compositions, as well as transfers in both directions between each
task-specific composition and the full SAM-40 composition (Full B), which
contains all three active tasks. All compositions contained the same
resting-state condition. For a transfer between two task-specific compositions,
the active task in the target composition was absent from the source
composition. In transfers involving Full B, the source and target compositions
overlapped in at least one active task: Full B included the task-specific
target when used as a source, and a task-specific source was included in Full B
when it was the target. These comparisons therefore test a change in the
mixture of task conditions, not generalization to an entirely unseen task.}

\rev{The cross-task folds were also subject-disjoint. For each fold, the source
composition from all but one participant formed the training set, whereas the
target composition from the held-out participant formed the test set.
Data-dependent preprocessing was performed separately within the selected
recordings of each participant and composition. Thus, no ICA decomposition
pooled source and target recordings, and the held-out target participant
contributed only to the ICA fitted on that participant's own unlabeled target
recordings. Every comparison therefore combines a change in participant
identity with a change in task composition. Only the transfers between two
task-specific compositions additionally evaluate an active task absent from
source training; none isolates the effect of task identity from the participant
shift. Keeping participants disjoint prevents participant-specific EEG
patterns from contributing to apparent transfer performance~\cite{kamrud2021vigilance,ke2021crosstask}.}

\textbf{Cross-dataset validation.}
\rev{To evaluate transfer between datasets, models were trained on all available data from a source composition and tested on a composition from the other dataset, in both transfer directions and without target-specific adaptation. Data-dependent preprocessing was performed separately for the two datasets: the mental attention-state dataset was cleaned using ICA fitted separately to each selected session, whereas SAM-40 was cleaned using ICA fitted separately to each participant's selected recordings. Consequently, source and target datasets were never pooled for ICA, and target recordings contributed only to a decomposition fitted on their own unlabeled acquisition grouping. SAM-40 was evaluated both as a complete composition and as task-specific compositions. Cross-dataset validation changes several factors simultaneously, including the participant population, experimental paradigm, recording setup, and condition definitions. It therefore measures transfer under their compound distribution shift rather than attributing a performance change to any single factor~\cite{lan2019domain}. The common sampling rate, harmonized binary labels, and shared 12-channel montage establish a compatible input and prediction space, but do not remove the substantive differences between the two datasets.}

\subsection*{\revs{Models}}

\rev{We evaluated five classifiers rather than relying on a single decoder to
test whether the conclusions are robust to a change in modeling assumptions.
This makes the decoder an explicit factor of the design: a difference between
validation protocols that is reproduced across models is more credibly
attributed to the data and evaluation protocol than to the idiosyncrasies of
one learner. The selected models span three families used in EEG decoding:
linear, tree-based, and deep-learning architectures. The deep networks are
included because they learn task-relevant representations directly from the
channel-by-time EEG windows, reducing reliance on a manually specified feature
set. All models were trained to predict whether a window belonged to a
high-attention or a low-attention state in the sense of
Table~\ref{tab:attention_label_mapping}.}

\rev{Hyperparameters were fixed and held constant across
datasets, task compositions, validation protocols and runs. No hyperparameter
search, early stopping or checkpoint selection was performed. Values followed library defaults or published reference
recipes, except for the departures specified below. A fixed random seed
governed fold construction, model initialization and network training. The
resolved configuration of each run was stored with its outputs.}

\textbf{Traditional machine learning.}
\rev{Both classical classifiers used the same fixed, label-free feature
representation. For each of the 12 channels, we computed log relative power in
the 1--4~Hz, 4--8~Hz, 8--13~Hz and 13--45~Hz bands, together with sample
entropy, mean, variance, skewness and kurtosis. Relative band power was
normalized by the channel total before log transformation. The nine features
from each channel formed one 108-dimensional feature vector per window.
This feature representation was applied unchanged to every dataset,
task composition, and validation protocol.}

\rev{Logistic regression provides a transparent linear baseline. It estimates
a linear decision boundary from the feature vectors, allowing the contribution
of individual features to be inspected. Its feature vectors were z-score
standardized using the mean and standard deviation estimated from the training
fold only.}

\rev{XGBoost~\cite{chen2016xgboost} represents tree-based ensemble learning,
allowing nonlinear relations and feature interactions to be modeled within the
same feature representation. It was also the decoder used in the originally
submitted version of this study, preserving a direct point of comparison with
those results. It received the same feature vectors without scaling. For both
classical models, inverse-frequency class weights were calculated from the
training fold only.}

\textbf{Deep learning.}
\rev{The remaining three classifiers operate directly on the prepared
channel-by-time EEG windows, rather than on the classical feature vectors. Each
input window had a $12 \times 640$ channel-by-time shape and was standardized
using statistics estimated from the training fold only.}

\rev{\textit{ShallowFBCSPNet}~\cite{schirrmeister2017deep} is a shallow
convolutional neural network designed to emulate the main operations of the
filter-bank common spatial pattern (FBCSP) pipeline within a trainable model.
It first applies temporal convolution and spatial filtering across channels,
then uses a squaring nonlinearity, mean pooling and a logarithmic activation to
form band-power-like features. Thus, its inductive bias reflects a conventional
EEG decoding approach, while its filters are optimized jointly with the
classifier.}

\rev{\textit{EEGNet}~\cite{lawhern2018eegnet} is a compact convolutional neural
network designed specifically for EEG. Its temporal convolutions learn filters
over time, depthwise convolutions learn spatial filters across electrodes for
each temporal feature map, and separable convolutions combine these learned
features efficiently. This factorized design substantially limits the number
of trainable parameters while retaining separate temporal and spatial feature
learning.}

\rev{\textit{EEG Conformer}~\cite{song2023conformer} combines a convolutional
front end with a Transformer encoder. The convolutional module forms local
temporal and spatial features and converts them into a sequence of embeddings.
Self-attention in the Transformer then models global dependencies between these
embeddings. It therefore complements the local receptive fields of convolution
with a mechanism for relating EEG patterns over longer time spans. To limit the
risk of overfitting in the smallest group-disjoint training folds, we reduced
the capacity of the Braindecode default configuration before evaluation and
then fixed it across all scenarios. Specifically, we reduced the number of
temporal convolution filters, which also determines the embedding width, from
40 to 16, the number of Transformer encoder layers from 6 to 2, and the number
of attention heads from 10 to 4. This retained the convolution--attention
architecture while reducing its parameter count for the $12 \times 640$ input
windows from 527,746 to 170,466.}

\rev{All three networks were trained from scratch, with no pretrained weights,
for 100 epochs with a batch size of 64. Their loss functions used
inverse-frequency class weights calculated from the training fold only, and no
resampling was applied.}

\section*{Results and analysis}

To characterize zero-shot generalization across the prespecified validation
strategies, we evaluated the binary contrast between high-attention and
low-attention states with five decoders.
Table~\ref{tab:scenario_metrics_logistic_regression} reports decoder-specific
performance across these strategies. It presents
participant-level balanced accuracy, macro F1, MCC, ROC-AUC, and per-class
recall for logistic regression, the interpretable reference decoder.
Corresponding tables for XGBoost, EEGNet, ShallowFBCSPNet, and EEGConformer are
provided in S1 Table, S2 Table, S3 Table and S4 Table.

\input{tables/scenario_metrics_logistic_regression}

\subsection*{Generalization under increasingly stringent validation strategies}

To examine generalization under increasingly stringent validation strategies,
we compared directed cross-task and cross-dataset transfer with the
baseline scenario for each target composition. These comparisons place
performance under progressively less similar training and test data in relation
to a target-specific reference. They therefore assess whether the binary
contrast was retained without target-specific adaptation as the participant,
task composition, or dataset differed between training and test data.

Fig~\ref{fig:transfer_matrix} summarizes the directed transfers between SAM-40 task compositions and
cross-dataset comparisons. Rows are grouped by source composition and decoder,
and columns denote target compositions. Each nonmissing matrix cell reports
balanced accuracy calculated separately from the held-out windows of each
target participant and then averaged across target participants for the given
decoder and direction. Off-diagonal cells represent zero-shot transfers without
target-specific adaptation, whereas the gray diagonal reports the
target-specific baseline estimated with five-fold stratified group validation
that keeps recording units disjoint between training and test folds. For the
summary columns, we first averaged
participant-level estimates across the five decoders for each source composition and target. We then
averaged these decoder-level transfer and baseline values across targets and
calculated their difference.
Negative values therefore indicate lower transfer performance than the
applicable target baselines, whereas dashes identify comparisons that are not
applicable.

Fig~\ref{fig:baseline_vs_cross_subject} compares baseline and
leave-one-subject-out cross-subject validation for the mental attention-state
dataset and the SAM-40 dataset. For each decoder and scenario, balanced accuracy was
calculated separately from the held-out windows of each target participant and
then averaged across target participants, and the colored lines show these
decoder-specific means. Black points show the mean of the five decoder-specific
estimates. Their bars show 95\% percentile bootstrap intervals for this mean within
each scenario, obtained by resampling target participants with replacement
10,000 times. Each draw uses the same participant indices for all five
decoders, with five participants in the mental attention-state dataset and
40 in SAM-40. The intervals describe variation across sampled participants
conditional on the fitted models, without refitting them or treating windows
or cross-validation folds as independent observations.

\input{figures/transfer_matrix}

For cross-task transfer, the mean decreases from the corresponding
target-specific baselines, averaged across all five decoders, were $9.6$,
$6.3$, $6.6$, and $3.1$ percentage points for models trained on the Stroop, Arithmetic, and
Mirror task compositions and the SAM-40 dataset, respectively. The mean
decrease across these source compositions was $6.4$ percentage points; this summary
includes both task-specific pairs and comparisons involving Full B, whose active
tasks overlap with those of the task-specific compositions. For cross-dataset
transfer, the decoder-averaged decreases for models trained on the Stroop,
Arithmetic, and Mirror task compositions, the SAM-40 dataset, and the mental
attention-state dataset were $13.1$, $10.4$, $13.1$, $11.0$, and $20.9$ percentage points,
respectively. The mean decrease was $13.7$ percentage points. Thus, zero-shot
transfer yielded lower performance in every eligible comparison, with a larger
reduction when training and test data came from
different datasets.

\input{figures/baseline_vs_cross_subject}

In the mental attention-state dataset, mean balanced accuracy was 65.5\% in the
baseline scenario and 65.4\% in cross-subject validation, a decrease of 0.1
percentage points. In the SAM-40 dataset, it decreased from 69.1\% to 67.1\%,
a difference of 2.0 percentage points. These descriptive differences do not
establish equivalence between baseline and cross-subject performance.

With two exceptions, all decoder-specific estimates fell within the
corresponding 95\% bootstrap interval. In the SAM-40 baseline scenario,
ShallowFBCSPNet reached 72.0\%, which was slightly above the interval upper
bound of 71.6\%. In cross-subject validation, EEGConformer reached 70.7\%,
which was above the upper bound of 69.7\%.

\subsection*{Performance consistency across decoders}

To assess whether performance patterns across validation strategies depended on
decoder choice, Fig~\ref{fig:scenario_accuracy_by_decoder} compares all
validation scenarios across the five decoders. Each colored point represents
the balanced accuracy of one decoder, averaged across target participants in a
given scenario. The black point and horizontal bar show the mean
across the five decoders and its paired 95\% percentile bootstrap confidence
interval (CI), respectively, estimated from 10,000 target-participant
resamples. The same resample was used for all decoders to preserve their
pairing.

\input{figures/scenario_accuracy_by_decoder}

Across the 27 scenarios, 42 of 135 decoder-specific estimates lay outside the
corresponding 95\% bootstrap interval. Logistic Regression lay outside the
interval in seven scenarios, with four points below and three above the
interval. XGBoost lay outside the interval in four scenarios, with one point
below and three above the interval. All of these departures occurred in
cross-dataset transfer settings.

EEGNet lay outside the interval in 12 scenarios, always below it.
ShallowFBCSPNet lay outside the interval in 10 scenarios, always above it.
These were the two decoders with the most frequent departures from the
interval for the mean across decoders. Such departures are descriptive and
do not constitute pairwise tests of decoder performance.

EEGConformer lay outside the interval in nine scenarios, with two points below
and seven above the interval. Overall, these patterns show that decoder choice
materially influenced balanced accuracy across validation scenarios, including
both within-dataset baselines and transfer settings.

\section*{Discussion}

Across the data compositions and validation protocols examined here, zero-shot
discriminative performance declined as the difference between the training and
test data increased. Mean performance decreased by 0.1 and 2.0 percentage points for unseen
participants in the mental attention-state dataset and SAM-40, respectively,
without establishing equivalence to baseline performance. Reductions relative
to the target-specific baselines were larger when the training and test
compositions also differed by task and larger still when they came from
different datasets. Because these protocols changed several factors together,
the reductions cannot be attributed to task composition or dataset alone. Thus, for these operationally defined high- versus
low-attention contrasts, technical compatibility and a common binary label space
were insufficient to preserve target-specific performance under increasingly
heterogeneous zero-shot transfer conditions.

\medskip

The overall degradation pattern was observed across the five decoders, but the
absolute level of balanced accuracy varied materially with decoder choice.
This finding makes the decoder an important factor in comparative EEG
evaluation rather than an interchangeable implementation detail. In
particular, conclusions about a validation scenario should distinguish a
pattern that is reproduced across model families from the performance of any
single predeclared model recipe.

\medskip

These results provide pre-adaptation baselines for assessing
source--target suitability in the studied data compositions. They indicate that
nominally aligned labels and partially harmonized acquisition specifications do
not by themselves establish that a source dataset will support useful zero-shot
transfer to a target dataset or task composition. Such baselines can therefore
inform source selection, which should be based on empirical, directed
source--target comparisons rather than on nominal label alignment alone. By
relating each transfer result to the corresponding target-specific baseline,
the analysis distinguishes preservation of the contrast under transfer from
the discriminability of that contrast in the target data itself. The resulting
references also make it possible to evaluate whether subsequent
transfer-learning, domain-adaptation, or personalization methods provide a
meaningful improvement over matched zero-shot performance rather than merely
achieving a favorable absolute score.

\medskip

Several limitations delimit the interpretation of these findings. First, the
binary labels distinguish protocol conditions whose EEG differences may
reflect task demands, stress, arousal, sensory input, or motor responses as
well as attention. In SAM-40, active cognitive tasks are contrasted with rest
while listening to music, so successful classification does not identify
which of these factors drives discrimination. Second, the mental
attention-state dataset contains only five participants, limiting population
generalization and the precision of participant-level bootstrap intervals
despite the many windows per participant. The use of only two technically
comparable datasets further limits the diversity of the evaluated conditions. The analysis also used the 12 shared
channels, a fixed parsimonious preparation pipeline, and predeclared model
recipes without hyperparameter or architecture selection to provide comparable
pre-adaptation benchmarks rather than to maximize performance in each
scenario. The reported figures consequently do not represent the maximum
attainable performance for an individual scenario.

\medskip

Future work could usefully optimize hyperparameters and architectural choices
within training folds, which would help distinguish limitations of the
predeclared recipes from those imposed by the data compositions and transfer
scenarios. It would also be informative to assess whether domain-adaptation
methods offer value relative to the same models without adaptation. Such
comparisons could provide a clearer indication of the extent to which
adaptation improves transfer-learning performance beyond the pre-adaptation
benchmarks reported here. To preserve the interpretability of these estimates,
future comparisons would benefit from retaining group-disjoint,
leakage-resistant splits and from avoiding the use of target test data for
model selection or parameter tuning.

% PLOS ONE ending sections
%   Acknowledgments -> References -> Supporting information captions
% (https://journals.plos.org/plosone/s/submission-guidelines)
%
% The "Funding" section that appeared in the submitted version has been removed.
% PLOS states: "Do not include funding sources in the Acknowledgments or anywhere
% else in the manuscript file. Funding information should only be entered in the
% financial disclosure section" of Editorial Manager. The funder, grant number
% BR27198099 and the Role of Funder statement are carried in
% submission_metadata.md instead.

\nolinenumbers

\bibliography{refs}

\section*{Supporting information}

\supportinginformationstyle


\input{tables/scenario_metrics_xgboost}
\input{tables/scenario_metrics_eegnet}
\input{tables/scenario_metrics_shallownet}
\input{tables/scenario_metrics_eegconformer}

\end{document}
